%! AP Statistics — Unit 1, Lesson 1.9 — Student Packet Cover Sheet
\documentclass[10pt]{article}
\usepackage{apstats-article}
\usepackage{apstats-boxes}

%=============================================================================
\begin{document}

% ── Full-bleed navy banner ────────────────────────────────────────────────────
\begin{tikzpicture}[remember picture, overlay]
  \fill[navy] (current page.north west)
    rectangle ([yshift=-0.9in]current page.north east);
\end{tikzpicture}
\vspace{-0.8in}

\begin{center}
  {\color{white}\textbf{\LARGE AP Statistics}}\\[4pt]
  {\color{skymid}\large\textbf{Unit 1: Exploring One-Variable Data}}\\[6pt]
  {\color{white}\textbf{\large Lesson 1.9 \quad Comparing Distributions of a Quantitative Variable}}
\end{center}

\vspace{0.22in}

\namedateperiod

\vspace{0.18in}

% ── LEARNING TARGETS ─────────────────────────────────────────────────────────
\begin{learningtargetbox}
\small
\begin{enumerate}[leftmargin=1.5em, itemsep=5pt]
  \item \ldots{} read parallel (side-by-side) boxplots and a back-to-back stemplot
        to identify the five-number summary and spread for each group.
  \item \ldots{} compare the \textbf{centers} of two distributions using explicit
        comparative language that names both groups and includes context.
  \item \ldots{} compare the \textbf{spreads} of two distributions, explaining what
        a larger IQR or range means in context.
  \item \ldots{} write a complete \textbf{AP-quality comparison} that addresses
        shape, outliers, center, and spread for both groups.
\end{enumerate}
\end{learningtargetbox}

\vspace{0.14in}

% ── TABLE OF CONTENTS ────────────────────────────────────────────────────────
\begin{tocbox}
\small
\renewcommand{\arraystretch}{1.5}
\begin{tabularx}{\linewidth}{@{} c l X r @{}}
\rowcolor{sky}
\textbf{\#} & \textbf{Component} & \textbf{Description} & \textbf{Score} \\
\midrule
1 & Warm-Up        & Read parallel boxplots; compare centers and spreads; AP language & \blank{1.2cm} \\
2 & Guided Notes   & Parallel boxplots, back-to-back stemplots, comparative statements & \blank{1.2cm} \\
3 & Group Activity & Compute summaries; interpret parallel boxplots; write AP comparison & \blank{1.2cm} \\
4 & Exit Ticket    & Read parallel boxplots; write a full AP comparative response       & \blank{1.2cm} \\
5 & Homework       & Read displays; fill-in stemplot; multi-part AP comparison practice  & \blank{1.2cm} \\
\midrule
  & \multicolumn{2}{r}{\textbf{Total}} & \blank{1.2cm} \\
\end{tabularx}
\end{tocbox}

\vspace{0.14in}

% ── KEEP IN MIND ─────────────────────────────────────────────────────────────
\begin{remindbox}
\small
In Lessons 1.7--1.8 you built and read boxplots for a single group. Lesson 1.9
places two or more boxplots on the \textit{same scale} so the distributions can
be compared directly.

\vspace{6pt}
\begin{tabularx}{\linewidth}{@{} >{\centering\arraybackslash\columncolor{goldbg}}p{0.06\linewidth}
                                    X @{}}
\textbf{\large!} &
\textbf{Comparison requires the same scale.}\enspace
Two boxplots on separate number lines cannot be meaningfully compared ---
always place them on one shared axis. \\[6pt]
\textbf{\large!} &
\textbf{AP scoring demands explicit comparative language.}\enspace
``Group A's median is 45'' earns no credit by itself. You must write: ``Group A's
median (45) is 8 points \textit{higher than} Group B's median (37).'' \\[6pt]
\textbf{\large!} &
\textbf{Always compare both center and spread.}\enspace
A difference in medians with a large overlap of IQR boxes may mean less than it
appears. Name both and explain what each tells you. \\[6pt]
\textbf{\large!} &
\textbf{Use context every time.}\enspace
Never write ``Group A'' without saying what that group is. Reference the variable
and the units in every comparative sentence.
\end{tabularx}
\end{remindbox}

\end{document}
